%-----------------------------------------------------------------------------------
%   PACKAGES 
%-----------------------------------------------------------------------------------

\documentclass[a4paper,11pt,abstract=true]{scrartcl}

\usepackage[margin=.9in]{geometry}


\usepackage{mlmodern}


\usepackage{amsmath,amsthm,amssymb}
\usepackage{mathtools,mathrsfs}
\usepackage{centernot}
\usepackage{wrapfig}

\usepackage[dvipsnames]{xcolor} % moved up: must be loaded before tikz and tcolorbox

\usepackage{tikz,tikz-cd,asymptote}
\usetikzlibrary{positioning, fit, calc}

\usepackage[most]{tcolorbox}
\usepackage{enumerate}
\usepackage[a]{esvect}
\usepackage{hyperref} % loaded last among the packages to avoid clashes


\makeatletter
\def\@seccntformat#1{\textcolor{Red}{\hspace*{0.2em}$\textbf{\S}$\hspace{-0.3em} \csname the#1\endcsname}\hspace{0.5em}}
\makeatother

\usepackage[headsepline]{scrlayer-scrpage}
\clearpairofpagestyles{}
\ihead{Vector Spaces}
\ohead{\pagemark}
\setkomafont{pageheadfoot}{\normalfont}
\setkomafont{pagenumber}{\normalfont}
\pagestyle{scrheadings}

\setlength{\parindent}{0cm}

%-----------------------------------------------------------------------------------
%   COMMANDS
%-----------------------------------------------------------------------------------

\DeclarePairedDelimiter{\floor}{\lfloor}{\rfloor}

% Bold, coloured run-in heading used for definitions, theorems, examples, etc.


\newcommand{\nonadj}{\,\neg\textsf{adj}\,}
\newcommand{\adj}{\,\textsf{adj}\,}
\newcommand{\RR}{\mathbb{R}}
\newcommand{\ZZ}{\mathbb{Z}}
\newcommand{\NN}{\mathbb{N}}
\newcommand{\CC}{\mathbb{C}}
\newcommand{\QQ}{\mathbb{Q}}
\newcommand{\sgn}{\text{sgn\hspace{0.5mm}}}
\newcommand{\osc}{\text{osc\hspace{0.5mm}}}
\newcommand{\foralls}{\hspace{2mm}\forall}
\newcommand{\rimply}{\Rightarrow}
\newcommand{\limply}{\Leftarrow}
\newcommand{\underdash}{\rule{0.25cm}{0.1mm}\hspace{0.5mm}}
\newcommand{\undersmalldash}{\rule{0.18cm}{0.1mm}\hspace{0.3mm}}
\newcommand{\s}{\square}
\newcommand{\xr}{\xrightarrow}

\newcommand{\genstirlingI}[3]{%
  \genfrac{\lbrack}{\rbrack}{0pt}{#1}{#2}{#3}%
}
\newcommand{\genstirlingII}[3]{%
  \genfrac{\{}{\}}{0pt}{#1}{#2}{#3}%
}
\newcommand{\genlah}[3]{%
  \genfrac{\lfloor}{\rfloor}{0pt}{#1}{#2}{#3}%
}
\newcommand{\stirlingI}[2]{\genstirlingI{}{#1}{#2}}
\newcommand{\stirlingII}[2]{\genstirlingII{}{#1}{#2}}
\newcommand{\lah}[2]{\genlah{}{#1}{#2}}

%-----------------------------------------------------------------------------------
%   ENVIROMENTS
%-----------------------------------------------------------------------------------

\newtheoremstyle{dotlessP}{}{}{}{}{\color{RedViolet}\bfseries}{.}{ }{}
\theoremstyle{dotlessP}
\newtheorem{problem}{Problem}
\newtheoremstyle{dotlessS}{}{}{}{}{\color{ForestGreen}\bfseries}{.}{ }{}
\theoremstyle{dotlessS}
\newtheorem*{solution}{Solution}

\theoremstyle{remark}
\newtheorem*{remark}{Remark}
\theoremstyle{definition}
\newtheorem*{theorem}{Theorem}
\newtheorem*{lemma}{Lemma}

\tcbset{
  base/.style={
    arc=0mm, 
    bottomtitle=0.5mm,
    boxrule=0mm,
    colbacktitle=black!10!white, 
    coltitle=black, 
    fonttitle=\bfseries, 
    left=2.5mm,
    leftrule=1mm,
    right=3.5mm,
    title={#1},
    toptitle=0.75mm, 
  }
}

\definecolor{brandblue}{rgb}{0.34, 0.7, 1}

\newtcolorbox{mainbox}[1]{
  enhanced jigsaw,
  breakable,
  colframe=brandblue, 
  base={#1}
}

\newtcolorbox{subbox}[1]{
  colframe=black!30!white,
  base={#1}
}

\newtcolorbox{definition_box}{
enhanced jigsaw,
breakable,
pad at break*=1mm,
colback=lightgray!12!white,
colframe=RedViolet!45!black,
colbacktitle=Gray!38!white,
before={\vspace{4mm}},
after={\vspace{2mm}},
sharp corners
}


\newtcolorbox{problem_box}{
enhanced jigsaw,
breakable,
pad at break*=1mm,
colback=white!5!white,
colframe=gray!75!black,
colbacktitle=white!28!white,
sharp corners,
before={\vspace{6mm}},
after={\vspace{2mm}}
}

\newtcolorbox{problem_box_1}{
enhanced jigsaw,
breakable,
pad at break*=1mm,
colback=SeaGreen!5!white,
colframe=SeaGreen!75!black,
colbacktitle=Black!28!white,
sharp corners,
before={\vspace{6mm}},
after={\vspace{2mm}}
}

\newtcolorbox{problem_box_title}[1]{
enhanced jigsaw,
breakable,
pad at break*=1mm,
colback=Apricot!5!white,
colframe=Melon!75!black,
sharp corners,
before={\vspace{2mm}},
after={\vspace{2mm}},
title={#1}, 
fonttitle=\bfseries
}

\newtcolorbox{claim_box}{
enhanced jigsaw,
breakable,
pad at break*=1mm,
colback=white,
colframe=white,
borderline west={2pt}{0pt}{white},
sharp corners,
before={\vspace{1mm}},
after={\vspace{0.8mm}}
}

\newtcolorbox{example_box}{
enhanced jigsaw,
breakable,
pad at break*=1mm,
colback=lightgray!12!white,
colframe=lightgray!15!white,
borderline west={2pt}{0pt}{black},
sharp corners,
before={\vspace{2mm}},
after={\vspace{2mm}}
}

\newtcolorbox{remark_box}{
enhanced jigsaw,
breakable,
pad at break*=1mm,
colback=Apricot!5!white,
colframe=Melon!75!black,
sharp corners,
before={\vspace{2mm}},
after={\vspace{2mm}}
}


\begin{document}

%-----------------------------------------------------------------------------------
%   TOP MATTERS
%-----------------------------------------------------------------------------------


\title{Vector Spaces}
\author{Pratyay Sen}
\date{\today}

\maketitle

\tableofcontents

\clearpage


%-----------------------------------------------------------------------------------
%   MAIN BODY
%-----------------------------------------------------------------------------------

\section{Fields}

\textbf{Definition} A \emph{field} is a set $\mathbb{F}$ equipped with two operations,
$+$ and $*$
\[
+:\mathbb{F}\times \mathbb{F}\longrightarrow \mathbb{F}
\]
\[
*:\mathbb{F}\times \mathbb{F}\longrightarrow \mathbb{F}
\]
such that $0\neq1$ and, for all $a,b,c\in \mathbb{F}$, the following axioms hold:
\begin{enumerate}
    \item (Closure under $+$) $a+b\in \mathbb{F}$;
    \item (Closure under $*$) $a*b\in \mathbb{F}$;
    \item (Associativity of $+$) $(a+b)+c=a+(b+c)$;
    \item (Associativity of $*$) $(a*b)*c=a*(b*c)$;
    \item (Commutativity of $+$) $a+b=b+a$;
    \item (Commutativity of $*$) $a*b=b*a$;
    \item (Additive identity) there exists $0\in \mathbb{F}$ such that
    $a+0=a$;
    \item (Multiplicative identity) there exists $1\in \mathbb{F}$ such that
    $a*1=a$;
    \item (Additive inverses) for every $a\in \mathbb{F}$, there exists $-a\in \mathbb{F}$
    such that $a+(-a)=0$;
    \item (Multiplicative inverses) for every $a\in \mathbb{F}$ with $a\neq0$,
    there exists $a^{-1}\in \mathbb{F}$ such that $a*a^{-1}=1$;
    \item (Distributivity) $a*(b+c)=a*b+a*c$.
\end{enumerate}

\begin{remark_box}
\textbf{A field is a substructure of a ring:} A field is a commutative
ring with identity in which every nonzero element has an inverse over the operation $*$.
For example, over addition and multiplication, $\mathbb{Z}$ is a ring but not a field, since $2$ has no
multiplicative inverse in $\mathbb{Z}$, whereas $\mathbb{Q}$, $\mathbb{R}$,
and $\mathbb{C}$ are fields.
\end{remark_box}

\textbf{Definition} If $\mathbb{F}$ is a field and $n$ is a positive integer, then
\[
\mathbb{F}^n=\underbrace{\mathbb{F}\times \mathbb{F}\times\cdots\times \mathbb{F}}_{n\text{ times}}
=\{(x_1,x_2,\ldots,x_n):x_i\in \mathbb{F}\}.
\]
Thus, $\mathbb{F}^n$ is the set of all ordered $n$-tuples of elements of $\mathbb{F}$. $+$
and $*$ are defined componentwise:
\[
(x_1,\ldots,x_n)+(y_1,\ldots,y_n)
=(x_1+y_1,\ldots,x_n+y_n),
\]
\[
a*(x_1,\ldots,x_n)=(a*x_1,\ldots,a*x_n), \qquad a\in \mathbb{F}.
\]
With these operations, $\mathbb{F}^n$ is a vector space over $\mathbb{F}$.

\begin{definition_box}
\textbf{Definition ($\mathbb{F}$-linear subspace of $\mathbb{F}^n$)} A subset
$W\subseteq \mathbb{F}^n$ is an \emph{$\mathbb{F}$-linear subspace of $\mathbb{F}^n$} if it contains the
zero vector and is closed under vector addition and multiplication by scalars
from $\mathbb{F}$. Explicitly:
\begin{enumerate}
    \item $(0,\ldots,0)\in W$;
    \item if $u,v\in W$, then $u+v\in W$;
    \item if $a\in \mathbb{F}$ and $u\in W$, then $a*u\in W$.
\end{enumerate}
Equivalently, a nonempty subset $W\subseteq \mathbb{F}^n$ is an $\mathbb{F}$-linear subspace if
\[
a*u+b*v\in W
\]
for every $u,v\in W$ and $a,b\in \mathbb{F}$. Thus, $W$ is closed under all linear
combinations of its elements. This is the special case $V=\mathbb{F}^n$ of the general
definition of a subspace of a vector space $V$.
\end{definition_box}

\textbf{Definition} The notation $\mathbb{F}^\infty$ denotes the set of all infinite
sequences of elements of $\mathbb{F}$:
\[
\mathbb{F}^\infty=\{(x_1,x_2,x_3,\ldots):x_i\in \mathbb{F}\text{ for every }i\in\mathbb{N}\}.
\]
$+$ and $*$ are defined componentwise:
\[
(x_1,x_2,\ldots)+(y_1,y_2,\ldots)
=(x_1+y_1,x_2+y_2,\ldots),
\]
\[
a*(x_1,x_2,\ldots)=(a*x_1,a*x_2,\ldots), \qquad a\in \mathbb{F}.
\]
With these operations, $\mathbb{F}^\infty$ is a vector space over $\mathbb{F}$. Here, sequences
may have infinitely many nonzero entries.

\section{Vector Spaces}

\textbf{Definition} Let $\mathbb{F}$ be a field. A \emph{vector space over $\mathbb{F}$} is a
set $V$ equipped with two binary operations $+$ and $*$
\[
+:V\times V\longrightarrow V
\]
\[
*:\mathbb{F}\times V\longrightarrow V
\]
such that, for all $u,v,w\in V$ and $a,b\in \mathbb{F}$:
\begin{enumerate}
    \item (Commutativity of $+$) $u+v=v+u$;
    \item (Associativity of $+$) $(u+v)+w=u+(v+w)$;
    \item (Identity over $+$) there exists $0_V\in V$ such that $v+0_V=v$;
    \item (Inverse over $+$) for each $v\in V$, there exists $-v\in V$ such that $v+(-v)=0_V$;
    \item (Distributivity of scalar $*$ over vector $+$) $a*(u+v)=a*u+a*v$;
    \item (Distributivity of scalar $+$ over scalar $*$ on vectors) $(a+b)*v=a*v+b*v$;
    \item (Associativity of $*$) $a*(b*v)=(a*b)*v$;
    \item (Identity over $*$) $1*v=v$, where $1$ is the identity of $\mathbb{F}$ over $*$.
\end{enumerate}

\textbf{Note:} As one might have noticed by now, I am deliberately avoiding using the term ``addition'' for $+$ and ``multiplication'' for $*$, because there are other functions that satisfy the field axioms and vector space axioms, that are not addition and multiplication. One such example is given below:

\begin{example_box}
\textbf{Example of nonstandard operations:} On the set $\mathbb{R}$, define
\[
x\oplus y=x+y+1
\qquad\text{and}\qquad
x\odot y=xy+x+y,
\]
where the operations on the right-hand sides are ordinary addition and
multiplication. These new operations are not the ordinary ones; for example,
$1\oplus1=3$ and $1\odot1=3$.

The set $\mathbb{R}$ is a field under $\oplus$ and $\odot$. Its identity over $\oplus$ is $-1$, its  identity over $\odot$ is $0$, the inverse over $\oplus$ of
$x$ is $-x-2$, and the inverse over $\odot$ of $x\neq-1$ is
\[
-\frac{x}{x+1}.
\]
To see why all the field axioms hold, consider the bijection
\[
T:\mathbb{R}\longrightarrow\mathbb{R},\qquad T(x)=x+1.
\]
It satisfies
\[
T(x\oplus y)=T(x)+T(y)
\qquad\text{and}\qquad
T(x\odot y)=T(x)T(y).
\]
Thus, $T$ converts the new operations into the ordinary field operations on
$\mathbb{R}$, so the usual field axioms are inherited.

This field is also a vector space over itself: take $V=\mathbb{R}$, use
$\oplus$ for vector addition, and define scalar multiplication by
\[
a*v=a\odot v=av+a+v.
\]
The vector-space axioms then follow from the field axioms. This example shows
that the symbols used for vector addition and scalar multiplication represent
abstract operations and need not be the usual arithmetic operations. However, now that the reader has understood that the ``addition'' and ``multiplication'' being talked about here need not necessarily be the conventional addition and multiplication, I shall use those terms from now on.
\end{example_box}

\subsection{Subspaces and affine spaces}
\textbf{Definition (subspace or linear subspace)} Let $V$ be a vector space over a field $\mathbb{F}$.
A subset $U\subseteq V$ is a \emph{subspace of $V$} if $U$ is itself a vector
space over $\mathbb{F}$ under the addition and scalar multiplication inherited from
$V$. Equivalently, $U$ is a subspace if:
\begin{enumerate}
    \item $0_V\in U$;
    \item $u+v\in U$ for every $u,v\in U$;
    \item $a*u\in U$ for every $a\in \mathbb{F}$ and $u\in U$.
\end{enumerate}
Equivalently, a nonempty subset $U\subseteq V$ is a subspace if
\[
a*u+b*v\in U
\]
for every $u,v\in U$ and $a,b\in \mathbb{F}$.

\begin{example_box}
\textbf{Trivial subspace:} The set containing only the zero
vector,
\[
\{0_V\},
\]
is a subspace of every vector space $V$ and is called the \emph{trivial
subspace} or \emph{zero subspace}. It is closed under vector addition and
scalar multiplication because
\[
0_V+0_V=0_V
\qquad\text{and}\qquad
a*0_V=0_V
\]
for every $a\in \mathbb{F}$. The subspaces $\{0_V\}$ and $V$ itself are sometimes
collectively called the trivial subspaces of $V$; all other subspaces are
called nontrivial proper subspaces.
\end{example_box}

\begin{definition_box}
\textbf{Definition (translate of a subspace)} Let $U$ be a subspace of $V$
and let $v\in V$. The \emph{translate} of $U$ by $v$ is
\[
v+U=\{v+u:u\in U\}.
\]
\end{definition_box}

\textbf{Definition (direction subspace and affine subspace)} A nonempty subset $A\subseteq V$ is an
\emph{affine subspace} of $V$ if there exist a linear subspace $U\subseteq V$
and a vector $v\in V$ such that
\[
A=v+U.
\]
The subspace $U$ is called the \emph{direction subspace} of $A$. Although the
vector $v$ is not unique, the direction subspace $U$ is unique. An affine
subspace looks like a shifted copy of $U$, but it need not contain $0_V$ and
therefore need not be a linear subspace. In fact,
\[
v+U\text{ is a subspace of }V
\quad\Longleftrightarrow\quad
v\in U,
\]
in which case $v+U=U$. Basically, 
every translate $A=v+U$ is an affine subspace with direction subspace $U$.  It is a linear
subspace exactly when $v\in U$, in which case $A=U$.

\textbf{Definition (coset)} Let $U$ be a subspace of $V$ and $v\in V$. The
set
\[
v+U=\{v+u:u\in U\}
\]
is called the \emph{coset of $U$ represented by $v$}. Thus, in a vector
space, cosets of linear subspaces are precisely affine subspaces. Two vectors
may represent the same coset; specifically,
\[
v+U=w+U
\quad\Longleftrightarrow\quad
v-w\in U.
\]
Consequently, any two cosets of the same subspace are either equal or
disjoint (prove it yourself!).

\textbf{Definition (quotient space)} Let $U$ be a subspace of a vector
space $V$ over $\mathbb{F}$. The \emph{quotient space of $V$ by $U$} is
\[
V/U=\{v+U:v\in V\},
\]
the set of all cosets of $U$ in $V$, equipped with the operations
\[
(v+U)+(w+U)=(v+w)+U,
\qquad
a*(v+U)=(a*v)+U,\quad a\in\mathbb{F}.
\]
Thus the elements of $V/U$ are entire cosets, not individual vectors of $V$.
Two vectors represent the same element of $V/U$ exactly when their
difference belongs to $U$.

These operations are \emph{well-defined}: their results do not depend on
which representatives we choose. Indeed, suppose $v+U=v'+U$ and
$w+U=w'+U$. Then $v-v'\in U$ and $w-w'\in U$. Since $U$ is a subspace,
\[
(v+w)-(v'+w')=(v-v')+(w-w')\in U,
\]
and, for every $a\in\mathbb{F}$,
\[
a*v-a*v'=a*(v-v')\in U.
\]
Hence replacing representatives leaves the resulting cosets unchanged.
The vector-space axioms follow from those of $V$. In particular, the zero
vector of $V/U$ is the coset $0_V+U=U$, and the additive inverse of $v+U$
is $(-v)+U$.

\begin{example_box}
\textbf{Example (quotient by the horizontal axis)} Take
\[
V=\mathbb{R}^2,
\qquad U=\{(t,0):t\in\mathbb{R}\}.
\]
For any $(a,b)\in V$,
\[
(a,b)+U=\{(t,b):t\in\mathbb{R}\}.
\]
Thus each element of $V/U$ is a horizontal line, determined entirely by its
height $b$. Adding two cosets adds their heights, and multiplying a coset by
a scalar multiplies its height. For example,
\[
\bigl((0,2)+U\bigr)+\bigl((0,3)+U\bigr)=(0,5)+U.
\]
The zero element of the quotient is the entire horizontal axis $U$.
\end{example_box}

\textbf{Note:} Please don't get discouraged upon seeing the number of definitions and propositions and theorems, the definitions are mostly just different ways of naming different stuff, and the propositions are just the immediate observable results. I am not saying that all propositions are trivial, but most of them are. 

\textbf{Examples.}
\begin{enumerate}
    \item Let $\mathbb{F}=\mathbb{R}$, $V=\mathbb{R}^3$, and
    \[
    U=\{(x,y,z)\in\mathbb{R}^3:x+y+z=0\}.
    \]
    Then $V$ is a vector space over $\mathbb{R}$, and $U$ is a subspace of
    $V$.

    \item Let $\mathbb{F}=\mathbb{C}$, $V=\mathbb{C}^2$, and
    \[
    U=\{(z,iz):z\in\mathbb{C}\}.
    \]
    Then $V$ is a vector space over $\mathbb{C}$, and $U$ is a subspace of
    $V$. Indeed, sums of vectors of the form $(z,iz)$ and complex scalar
    multiples of such vectors have the same form.

    \item Let $\mathbb{F}=\mathbb{R}$, let $V=\mathcal{P}_3(\mathbb{R})$ be the vector
    space of real polynomials of degree at most $3$, and let
    \[
    U=\{p\in\mathcal{P}_3(\mathbb{R}):p(0)=0\}.
    \]
    Then $U$ is a subspace of $V$, since the zero polynomial belongs to $U$
    and, whenever $p(0)=q(0)=0$, we have
    $(p+q)(0)=0$ and $(a*p)(0)=0$ for every $a\in\mathbb{R}$.
\end{enumerate}

\subsection{Sums and intersections}
\begin{definition_box}
\textbf{Sum of Subspaces} Let $U_1,\ldots,U_m$ be subspaces of
a vector space $V$ over $\mathbb{F}$. Their \emph{sum} is
\[
U_1+\cdots+U_m
=\{u_1+\cdots+u_m:u_i\in U_i\text{ for }i=1,\ldots,m\}.
\]
Thus, the sum consists of all vectors that can be written as a sum of one
vector from each subspace. In particular,
\[
U+W=\{u+w:u\in U,\ w\in W\}.
\]
\end{definition_box}

\textbf{Proposition:} The sum $U_1+\cdots+U_m$ is a subspace of $V$.

\textit{Proof:} Because each $U_i$ is a subspace, $0_V\in U_i$; hence
\[
0_V=0_V+\cdots+0_V\in U_1+\cdots+U_m.
\]
Now let
\[
x=u_1+\cdots+u_m
\quad\text{and}\quad
y=v_1+\cdots+v_m
\]
belong to $U_1+\cdots+U_m$, where $u_i,v_i\in U_i$. Then
\[
x+y=(u_1+v_1)+\cdots+(u_m+v_m).
\]
Since each $U_i$ is closed under addition, $u_i+v_i\in U_i$, and therefore
$x+y\in U_1+\cdots+U_m$. Finally, for any $a\in \mathbb{F}$,
\[
a*x=(a*u_1)+\cdots+(a*u_m).
\]
Since each $U_i$ is closed under scalar multiplication, $a*u_i\in U_i$, so
$a*x\in U_1+\cdots+U_m$. Thus the sum contains $0_V$ and is closed under
addition and scalar multiplication; hence it is a subspace of $V$.
\hfill$\square$

\textbf{Proposition:} The sum $U_1+\cdots+U_m$ is the smallest subspace of
$V$ containing $U_1,\ldots,U_m$.

\textit{Proof:} First, the sum contains each $U_j$. Indeed, if $u\in U_j$,
then
\[
u=0_V+\cdots+0_V+u+0_V+\cdots+0_V\in U_1+\cdots+U_m,
\]
where $u$ occupies the $j$th position and every other term is the zero vector
of the corresponding subspace. Hence
\[
U_j\subseteq U_1+\cdots+U_m
\qquad\text{for every }j=1,\ldots,m.
\]

Now let $W$ be any subspace of $V$ containing every $U_j$. If
$x\in U_1+\cdots+U_m$, then
\[
x=u_1+\cdots+u_m
\]
for some $u_j\in U_j$. Since $U_j\subseteq W$, every $u_j$ belongs to $W$.
Because $W$ is closed under addition, $x\in W$. Therefore
\[
U_1+\cdots+U_m\subseteq W.
\]
Thus, $U_1+\cdots+U_m$ contains all the subspaces $U_j$ and is contained in
every subspace that contains them. It is therefore the smallest such subspace.
\hfill$\square$

\textbf{Proposition (intersection of two subspaces):} If $U$ and $W$ are
subspaces of $V$, then $U\cap W$ is a subspace of $V$.
\textit{Proof:} Because $U$ and $W$ are subspaces, both contain the zero
vector. Hence
\[
0_V\in U\cap W.
\]
Now let $u,v\in U\cap W$. Then $u,v\in U$ and $u,v\in W$. Since both
subspaces are closed under addition,
\[
u+v\in U
\qquad\text{and}\qquad
u+v\in W,
\]
so $u+v\in U\cap W$. Finally, if $a\in \mathbb{F}$, then closure of $U$ and $W$ under
scalar multiplication gives
\[
a*u\in U
\qquad\text{and}\qquad
a*u\in W.
\]
Therefore $a*u\in U\cap W$. Thus $U\cap W$ contains $0_V$ and is closed
under addition and scalar multiplication, so it is a subspace of $V$.
\hfill$\square$

\textbf{Proposition (union of two subspaces):} If $U$ and $W$ are subspaces
of $V$, then $U\cup W$ is a subspace of $V$ if and only if
\[
U\subseteq W
\qquad\text{or}\qquad
W\subseteq U.
\]

\textit{Proof:} If $U\subseteq W$, then $U\cup W=W$, which is a subspace.
Similarly, if $W\subseteq U$, then $U\cup W=U$, which is a subspace.

Conversely, suppose that $U\cup W$ is a subspace. Assume neither subspace is
contained in the other. Then there exist
\[
u\in U\setminus W
\qquad\text{and}\qquad
w\in W\setminus U.
\]
Because $U\cup W$ is assumed to be a subspace and $u,w\in U\cup W$, it must
contain $u+w$. Thus $u+w$ belongs to either $U$ or $W$. If $u+w\in U$, then
\[
w=(u+w)-u\in U,
\]
contradicting $w\notin U$. If $u+w\in W$, then
\[
u=(u+w)-w\in W,
\]
contradicting $u\notin W$. Hence one of $U$ and $W$ must be contained in the
other. \hfill$\square$

\subsection{Spanning and linear independence}
\begin{definition_box}
\textbf{Definition (span)} Let $v_1,\ldots,v_m\in V$. The \emph{span} of
$v_1,\ldots,v_m$ is the set of all their linear combinations:
\[
\operatorname{span}(v_1,\ldots,v_m)
=\{a_1*v_1+\cdots+a_m*v_m:a_1,\ldots,a_m\in \mathbb{F}\}.
\]
More generally, if $S\subseteq V$, then
\[
\operatorname{span}(S)
=\left\{a_1*v_1+\cdots+a_m*v_m:
m\in\mathbb{N},\ v_i\in S,\ a_i\in \mathbb{F}\right\}.
\]
\end{definition_box}

Thus, $\operatorname{span}(S)$ is the smallest subspace of $V$ containing
$S$. By convention, $\operatorname{span}(\varnothing)=\{0_V\}$.

\begin{definition_box}
\textbf{Definition (linear independence)} A list of vectors
$v_1,\ldots,v_m\in V$ is \emph{linearly independent} over $\mathbb{F}$ if no vector
in the list belongs to the span of the others; that is,
\[
v_j\notin\operatorname{span}(v_1,\ldots,v_{j-1},v_{j+1},\ldots,v_m)
\]
for every $j=1,\ldots,m$. Equivalently, the list is linearly independent if
\[
a_1*v_1+\cdots+a_m*v_m=0_V
\]
implies
\[
a_1=\cdots=a_m=0.
\]
In other words, the zero vector can be written as a linear combination of
$v_1,\ldots,v_m$ only by taking every coefficient to be zero. Another
equivalent span criterion is
\[
v_j\notin\operatorname{span}(v_1,\ldots,v_{j-1})
\qquad\text{for every }j=1,\ldots,m,
\]
where $\operatorname{span}(\varnothing)=\{0_V\}$.
\end{definition_box}

\begin{definition_box}
\textbf{Definition (linear dependence)} A list of vectors
$v_1,\ldots,v_m\in V$ is \emph{linearly dependent} over $\mathbb{F}$ if it is not
linearly independent; that is, if there exist scalars
$a_1,\ldots,a_m\in \mathbb{F}$, not all zero, such that
\[
a_1*v_1+\cdots+a_m*v_m=0_V.
\]
Such an equation is called a \emph{nontrivial linear relation}. Equivalently,
a finite list is linearly dependent if one of its vectors lies in the span of
the others.
\end{definition_box}


\textbf{Definition (generating set)} A subset $S\subseteq V$ is called a
\emph{generating set} or \emph{spanning set} for $V$ if
\[
\operatorname{span}(S)=V.
\]
In this case, we say that $S$ \emph{generates} or \emph{spans} $V$. This
means that every vector $v\in V$ can be written as a finite linear
combination of elements of $S$.

More generally, if $S\subseteq V$, then $S$ generates the subspace
$\operatorname{span}(S)$, even when it does not generate all of $V$. A
generating set need not be linearly independent; a basis is precisely a
linearly independent generating set.



\textbf{Definition (finitely generated vector space)} A vector space $V$
over $\mathbb{F}$ is \emph{finitely generated} if it has a finite generating set. In
other words, there exist vectors $v_1,\ldots,v_n\in V$ such that
\[
V=\operatorname{span}(v_1,\ldots,v_n).
\]
For vector spaces, being finitely generated is equivalent to being
finite-dimensional: from any finite generating set, redundant vectors can be
removed until a finite basis remains. Conversely, every finite basis is a
finite generating set.


\subsection{Bases and dimension}
\begin{definition_box}
\textbf{Definition (basis)} A list of vectors
$v_1,\ldots,v_m\in V$ is a \emph{basis} of $V$ if:
\begin{enumerate}
    \item $v_1,\ldots,v_m$ are linearly independent; and
    \item $\operatorname{span}(v_1,\ldots,v_m)=V$.
\end{enumerate}
Equivalently, $v_1,\ldots,v_m$ form a basis of $V$ if every vector $v\in V$
can be written uniquely as
\[
v=a_1*v_1+\cdots+a_m*v_m
\]
for scalars $a_1,\ldots,a_m\in \mathbb{F}$. The scalars $a_1,\ldots,a_m$ are called
the \emph{coordinates} of $v$ with respect to this basis.
\end{definition_box}

\textbf{Theorem} If a vector space $V$ over $\mathbb{F}$ is generated by $n$
vectors, then $V$ has a basis containing at most $n$ vectors.

\textit{Proof:} Suppose
\[
V=\operatorname{span}(v_1,\ldots,v_n).
\]
If $(v_1,\ldots,v_n)$ is linearly independent, then it is already a basis of
$V$, and the result follows.

If the list is linearly dependent, then one of its vectors belongs to the
span of the others. After relabeling if necessary, suppose
\[
v_n\in\operatorname{span}(v_1,\ldots,v_{n-1}).
\]
Removing $v_n$ does not change the span, so
\[
V=\operatorname{span}(v_1,\ldots,v_{n-1}).
\]
If this shorter list is still linearly dependent, remove another vector that
lies in the span of the remaining vectors. Each removal preserves the span
and decreases the length of the list by one.

Because the original list is finite, the process eventually stops with a
linearly independent list $(u_1,\ldots,u_m)$ that still spans $V$. Hence this
list is a basis of $V$, and because it was obtained by deleting vectors from
the original list,
\[
m\leq n.
\]
Thus $V$ has a basis of at most $n$ vectors. \hfill$\square$

\textbf{Definition (basis expansion)} Let
$B=(v_1,\ldots,v_n)$ be a basis of $V$. For each $v\in V$, the unique
expression
\[
v=a_1*v_1+\cdots+a_n*v_n,
\qquad a_1,\ldots,a_n\in \mathbb{F},
\]
is called the \emph{basis expansion of $v$ with respect to $B$}. The scalars
$a_1,\ldots,a_n$ are the coordinates of $v$, and the tuple
\[
{[v]}_B=(a_1,\ldots,a_n)\in \mathbb{F}^n
\]
is called the \emph{coordinate vector of $v$ relative to $B$}.

If $B={(v_i)}_{i\in I}$ is an infinite basis, then the basis expansion of $v$
has the form
\[
v=\sum_{i\in I}a_i*v_i,
\]
where only finitely many coefficients $a_i$ are nonzero. Thus every basis
expansion is a finite linear combination, even when the basis is infinite.

\textbf{Definition (dimension)} The \emph{dimension} of a vector space $V$
over $\mathbb{F}$, denoted by $\dim_{\mathbb{F}} V$ or simply $\dim V$, is the number of vectors
in any basis of $V$. All bases of a vector space have the same number of
elements, so the dimension is well-defined. If
\[
(v_1,\ldots,v_n)
\]
is a basis of $V$, then $\dim_{\mathbb{F}} V=n$. If no finite basis exists, then $V$ is
called \emph{infinite-dimensional}. The zero vector space has the empty basis,
and hence
\[
\dim_{\mathbb{F}}\{0_V\}=0.
\]

\begin{mainbox}{All Finite Bases have the Same Length}
If
$(v_1,\ldots,v_n)$ and $(w_1,\ldots,w_m)$ are two finite bases of $V$, then
$m=n$.
\end{mainbox}

\textit{Proof:} We first show that $m\leq n$. Because
$(v_1,\ldots,v_n)$ spans $V$, we can write
\[
w_1=a_1*v_1+\cdots+a_n*v_n.
\]
At least one coefficient is nonzero. After relabeling the $v_i$ if necessary,
suppose $a_1\neq0$. We can then solve for $v_1$:
\[
v_1=a_1^{-1}*w_1-a_1^{-1}*a_2*v_2-\cdots-a_1^{-1}*a_n*v_n.
\]
Therefore
\[
(w_1,v_2,\ldots,v_n)
\]
still spans $V$.

Now write $w_2$ as a linear combination of this new spanning list. The
coefficient of at least one of $v_2,\ldots,v_n$ must be nonzero; otherwise
$w_2$ would belong to $\operatorname{span}(w_1)$, contradicting the linear
independence of $(w_1,w_2)$. Relabeling the remaining $v_i$ if necessary, we
can solve for $v_2$ and replace it by $w_2$. Hence
\[
(w_1,w_2,v_3,\ldots,v_n)
\]
still spans $V$.

Continuing in this way, after $k$ steps we obtain a spanning list
\[
(w_1,\ldots,w_k,v_{k+1},\ldots,v_n).
\]
This process cannot continue beyond $k=n$. If $m>n$, then after replacing all
the $v_i$, the list $(w_1,\ldots,w_n)$ would span $V$. Consequently,
$w_{n+1}$ would lie in $\operatorname{span}(w_1,\ldots,w_n)$, contradicting
the linear independence of $(w_1,\ldots,w_m)$. Thus $m\leq n$.

Interchanging the roles of the two bases gives $n\leq m$. Therefore
$m=n$. \hfill$\square$

\begin{mainbox}{Size and Extension of Linearly Independent Sets} Let $V$
be a vector space over $\mathbb{F}$ with $\dim_{\mathbb{F}} V=n$. If $L\subseteq V$ is linearly
independent, then
\[
|L|\leq n.
\]
Equality holds if and only if $L$ is a basis of $V$. If $|L|<n$, then $L$
can be extended to a basis of $V$.
\end{mainbox}

\textit{Proof:} Fix a basis $B=(b_1,\ldots,b_n)$ of $V$. The replacement
argument used above shows that an independent list cannot contain more
vectors than a spanning list. Since $B$ spans $V$, a linearly independent set
$L$ cannot contain $n+1$ elements. Hence $L$ is finite and $|L|\leq n$.

Write $L=\{u_1,\ldots,u_m\}$, where $m\leq n$. Beginning with the spanning
list $(b_1,\ldots,b_n)$, write $u_1$ as a linear combination of its vectors.
At least one coefficient is nonzero, so we can solve for the corresponding
$b_i$ and replace it by $u_1$ without changing the span.

Suppose that $u_1,\ldots,u_{k-1}$ have already replaced $k-1$ vectors of
$B$. Express $u_k$ as a linear combination of the resulting spanning list.
The coefficient of at least one remaining vector from $B$ must be nonzero;
otherwise $u_k$ would lie in
$\operatorname{span}(u_1,\ldots,u_{k-1})$, contradicting the linear
independence of $L$. We may therefore solve for that basis vector and replace
it by $u_k$ while preserving the span. Continuing in this way gives a
spanning list of the form
\[
(u_1,\ldots,u_m,c_{m+1},\ldots,c_n),
\]
where the vectors $c_{m+1},\ldots,c_n$ were chosen from $B$. This list has
$n$ elements. It must be linearly independent: otherwise one of its vectors
could be removed while preserving the span, producing a spanning list of
$n-1$ vectors. This is impossible because the replacement argument shows
that the independent list $B$, which has $n$ vectors, cannot be longer than
a spanning list. Thus the displayed list is a basis extending $L$.

If $m=n$, there are no additional vectors $c_i$, so $L$ itself is a basis.
Conversely, if $L$ is a basis, then the theorem that all finite bases have the
same length gives $|L|=n$. If $m<n$, the vectors
$c_{m+1},\ldots,c_n$ extend $L$ to a basis. \hfill$\square$

\begin{mainbox}{Bases of Subspaces} 
Every subspace of a
finite-dimensional vector space has a finite basis.
\end{mainbox}

\textit{Proof:} Let $U$ be a subspace of a finite-dimensional vector space
$V$, and write $\dim_{\mathbb{F}} V=n$. Every linearly independent subset of $U$ is also
a linearly independent subset of $V$, so the preceding theorem shows that it
has at most $n$ elements. Therefore, among all linearly independent subsets
of $U$, choose one, say $L$, having the greatest possible number of elements.

We claim that $L$ spans $U$. If it did not, there would exist
\[
u\in U\setminus\operatorname{span}(L).
\]
Then $L\cup\{u\}$ would be linearly independent: in any relation involving
$u$, a nonzero coefficient of $u$ would imply
$u\in\operatorname{span}(L)$, while a zero coefficient would reduce the
relation to one among the elements of $L$. Thus all coefficients must be
zero. But $L\cup\{u\}$ has more elements than $L$, contradicting the maximal
choice of $L$. Hence $L$ spans $U$. Since $L$ is both linearly independent
and spanning, it is a basis of $U$. \hfill$\square$

\begin{mainbox}{Size and Reduction of Spanning Sets} 
Let $V$ be a vector
space over $\mathbb{F}$ with $\dim_{\mathbb{F}} V=n$. If $S\subseteq V$ spans $V$, then
\[
|S|\geq n.
\]
Equality holds if and only if $S$ is a basis of $V$. If $|S|>n$, then $S$
contains a subset that is a basis of $V$.
\end{mainbox}

\textit{Proof:} Let $B=(b_1,\ldots,b_n)$ be a basis of $V$. Because $S$
spans $V$, each $b_i$ is a finite linear combination of elements of $S$.
For each $i$, choose a finite subset $S_i\subseteq S$ such that
\[
b_i\in\operatorname{span}(S_i),
\]
and set
\[
S_0=S_1\cup\cdots\cup S_n.
\]
The set $S_0$ is finite and spans every element of $B$; hence it spans $V$.
By repeatedly deleting any element of $S_0$ that lies in the span of the
others, we obtain a linearly independent spanning subset $B'\subseteq S_0$.
Thus $B'$ is a basis of $V$, and
\[
B'\subseteq S_0\subseteq S.
\]
Because every basis of $V$ has $n$ elements, $|B'|=n$, and therefore
$|S|\geq n$.

If $|S|=n$, then the inclusion $B'\subseteq S$ and the equality
$|B'|=|S|=n$ imply $B'=S$, so $S$ is a basis. Conversely, every basis spans
$V$ and has exactly $n$ elements. Finally, if $|S|>n$, the subset
$B'\subseteq S$ constructed above is a basis of $V$. \hfill$\square$

\begin{definition_box}
\textbf{Definition (standard basis of $\mathbb{F}^n$).} For $j=1,\ldots,n$, let
\[
e_j=(0,\ldots,0,1,0,\ldots,0),
\]
where the $1$ occurs in the $j$th position. The list
\[
(e_1,\ldots,e_n)
\]
is called the \emph{standard basis} of $\mathbb{F}^n$. Every
$x=(x_1,\ldots,x_n)\in \mathbb{F}^n$ has the unique representation
\[
x=x_1*e_1+\cdots+x_n*e_n,
\]
and therefore $\dim_{\mathbb{F}} \mathbb{F}^n=n$.
\end{definition_box}

\begin{example_box}
\textbf{Other standard bases.}
\begin{enumerate}
    \item The standard basis of the polynomial space
    $\mathcal{P}_n(\mathbb{F})$ is
    \[
    (1,x,x^2,\ldots,x^n),
    \]
    so $\dim_{\mathbb{F}}\mathcal{P}_n(\mathbb{F})=n+1$.

    \item For the matrix space $\mathbb{F}^{m\times n}$, let $E_{ij}$ be the matrix
    whose $(i,j)$ entry is $1$ and whose other entries are $0$. The matrices
    \[
    E_{11},E_{12},\ldots,E_{mn}
    \]
    form its standard basis, so $\dim_{\mathbb{F}} \mathbb{F}^{m\times n}=mn$.
\end{enumerate}
\end{example_box}

\textbf{Proposition:} Any three polynomials in $\mathcal{P}_1(\mathbb{F})$ are
linearly dependent.

\textit{Proof:} The list $(1,x)$ is a basis of $\mathcal{P}_1(\mathbb{F})$, so
\[
\dim_{\mathbb{F}}\mathcal{P}_1(\mathbb{F})=2.
\]
Suppose $p_1,p_2,p_3\in\mathcal{P}_1(\mathbb{F})$. If $p_1$ and $p_2$ are linearly
dependent, then $(p_1,p_2,p_3)$ is also linearly dependent. Otherwise,
$(p_1,p_2)$ is a linearly independent list of two vectors in a
two-dimensional space and therefore forms a basis. Hence
\[
p_3=a*p_1+b*p_2
\]
for some $a,b\in \mathbb{F}$. Rearranging gives
\[
a*p_1+b*p_2-1*p_3=0.
\]
Because the coefficient of $p_3$ is $-1\neq0$, this is a nontrivial linear
relation. Thus $(p_1,p_2,p_3)$ is linearly dependent. \hfill$\square$

\textbf{Proposition (dependence of polynomials):} Let $n\geq0$ be an
integer. Any set of $n+2$ polynomials in $\mathcal{P}_n(\mathbb{F})$, the
space of polynomials in $x$ of degree at most $n$ together with the zero
polynomial, is linearly dependent over $\mathbb{F}$.

\textit{Proof:} The monomials
\[
1,x,x^2,\ldots,x^n
\]
form a basis of $\mathcal{P}_n(\mathbb{F})$, so this space has dimension
$n+1$. By the theorem on the size of linearly independent sets, every
linearly independent subset of this space has at most $n+1$ elements.
Since $n+2>n+1$, a set $\{p_1,\ldots,p_{n+2}\}$ of $n+2$ polynomials
cannot be linearly independent. Thus there exist scalars
$a_1,\ldots,a_{n+2}\in\mathbb{F}$, not all zero, such that
\[
a_1*p_1+\cdots+a_{n+2}*p_{n+2}=0,
\]
where $0$ denotes the zero polynomial. Hence the set is linearly dependent.
\hfill$\square$

\begin{mainbox}{Dimension of a Quotient Space}
Let $V$ be a
finite-dimensional vector space over $\mathbb{F}$ and let $U$ be a subspace
of $V$. Then
\[
\dim(V/U)+\dim U=\dim V.
\]
\end{mainbox}

\textit{Proof:} Write $\dim U=k$ and $\dim V=n$. Choose a basis
$u_1,\ldots,u_k$ of $U$. By the basis-extension theorem,
extend it to a basis
\[
u_1,\ldots,u_k,v_1,\ldots,v_{n-k}
\]
of $V$. We claim that the cosets
\[
v_1+U,\ldots,v_{n-k}+U
\]
form a basis of $V/U$.

First, they span $V/U$. Every $v\in V$ has an expression
\[
v=\sum_{i=1}^{k}a_i u_i+\sum_{j=1}^{n-k}b_j v_j,
\qquad a_i,b_j\in\mathbb{F}.
\]
The first sum belongs to $U$, so in the quotient it contributes the zero
coset. Hence
\[
v+U=\sum_{j=1}^{n-k}b_j*(v_j+U).
\]

Next, suppose a linear combination of these cosets is zero in $V/U$:
\[
\sum_{j=1}^{n-k}c_j*(v_j+U)=U.
\]
This means $\sum_{j=1}^{n-k}c_j v_j\in U$, so there are scalars
$d_1,\ldots,d_k\in\mathbb{F}$ such that
\[
\sum_{j=1}^{n-k}c_j v_j-\sum_{i=1}^{k}d_i u_i=0_V.
\]
Since the extended basis of $V$ is linearly independent, all coefficients
are zero. In particular, every $c_j=0$, proving independence in $V/U$.

Thus $\dim(V/U)=n-k$, and therefore
\[
\dim(V/U)+\dim U=(n-k)+k=n=\dim V.
\]
Empty lists and empty sums are allowed here, so the proof also covers
$U=\{0_V\}$ and $U=V$. \hfill$\square$

\subsection{Row and column spaces}
\textbf{Row Space} Let $A\in \mathbb{F}^{m\times n}$ be a matrix whose
rows are $r_1,\ldots,r_m\in \mathbb{F}^n$. The \emph{row space} of $A$ is the subspace
of $\mathbb{F}^n$ spanned by its rows:
\[
\operatorname{Row}(A)
=\operatorname{span}(r_1,\ldots,r_m)
=\{a_1*r_1+\cdots+a_m*r_m:a_1,\ldots,a_m\in \mathbb{F}\}.
\]

\textbf{Column Space} Let the columns of
$A\in \mathbb{F}^{m\times n}$ be $c_1,\ldots,c_n\in \mathbb{F}^m$. The \emph{column space} of
$A$ is the subspace of $\mathbb{F}^m$ spanned by its columns:
\[
\operatorname{Col}(A)
=\operatorname{span}(c_1,\ldots,c_n)
=\{a_1*c_1+\cdots+a_n*c_n:a_1,\ldots,a_n\in \mathbb{F}\}.
\]
Thus, for an $m\times n$ matrix, the row space is contained in $\mathbb{F}^n$, whereas
the column space is contained in $\mathbb{F}^m$.

\subsection{Direct sums and complements}
\textbf{Internal Direct Sum} Let $U_1,\ldots,U_m$ be subspaces
of $V$. Their sum is called an \emph{internal direct sum} if every vector in
$U_1+\cdots+U_m$ has a unique representation of the form
\[
u_1+\cdots+u_m,
\qquad u_i\in U_i.
\]
In this case, we write
\[
U_1\oplus\cdots\oplus U_m
\]
instead of $U_1+\cdots+U_m$. Thus,
\[
V=U_1\oplus\cdots\oplus U_m
\]
means that every $v\in V$ can be written uniquely as
$v=u_1+\cdots+u_m$ with $u_i\in U_i$.

Equivalently, the sum is direct if
\[
u_1+\cdots+u_m=0_V,
\qquad u_i\in U_i,
\]
implies $u_1=\cdots=u_m=0_V$. Another equivalent condition is
\[
U_j\cap\sum_{i\neq j}U_i=\{0_V\}
\qquad\text{for every }j.
\]
For two subspaces, this simplifies to
\[
U+W=U\oplus W
\quad\Longleftrightarrow\quad
U\cap W=\{0_V\}.
\]

\textbf{Complementary Subspace} Let $U$ be a subspace of $V$.
A subspace $W\subseteq V$ is called a \emph{complementary subspace} or
\emph{complement} of $U$ in $V$ if
\[
V=U\oplus W.
\]
Equivalently, $W$ is a complement of $U$ if
\[
U+W=V
\qquad\text{and}\qquad
U\cap W=\{0_V\}.
\]
Thus every $v\in V$ has a unique decomposition
\[
v=u+w,
\qquad u\in U,\quad w\in W.
\]
A complementary subspace need not be unique. If $V$ is finite-dimensional,
then every subspace has a complement, and any complement $W$ satisfies
\[
\dim V=\dim U+\dim W.
\]

\textbf{Theorem (existence of complementary subspaces)} Let $V$ be
finite-dimensional and let $U$ be a subspace of $V$. Then there exists a
subspace $W\subseteq V$ such that
\[
V=U\oplus W.
\]

\textit{Proof:} Choose a basis $B_U$ of $U$. By the basis-extension theorem,
$B_U$ can be extended to a basis of $V$. Thus there is a set $C\subseteq V$
such that
\[
B_U\cup C
\]
is a basis of $V$. Define
\[
W=\operatorname{span}(C).
\]

Because $B_U\cup C$ spans $V$, every $v\in V$ can be written as a finite
linear combination of elements of $B_U\cup C$. Grouping the terms from
$B_U$ and the terms from $C$ gives
\[
v=u+w
\]
for some $u\in U$ and $w\in W$. Hence $U+W=V$.

It remains to show that $U\cap W=\{0_V\}$. Let $x\in U\cap W$. Because
$x\in U$, it is a linear combination of elements of $B_U$; because
$x\in W$, it is a linear combination of elements of $C$. Subtracting these
two expressions gives a linear combination of elements of $B_U\cup C$ equal
to $0_V$. Since $B_U\cup C$ is linearly independent, every coefficient must
be zero. Therefore $x=0_V$, and thus
\[
U\cap W=\{0_V\}.
\]
Consequently, $V=U\oplus W$, so $W$ is a complementary subspace of $U$.
\hfill$\square$

\begin{remark_box}
\textbf{Is every sum a direct sum?} Nope, this statement is not always true. Every direct sum is a sum, but a sum
is direct only when the decomposition of each vector is unique. For two
subspaces, the sum is direct precisely when $U\cap W=\{0_V\}$.
\end{remark_box}

\begin{example_box}
\textbf{Example.} In $V=\mathbb{R}^2$, let
\[
U=\operatorname{span}((1,0))
\qquad\text{and}\qquad
W=\mathbb{R}^2.
\]
Then
\[
U+W=\mathbb{R}^2,
\]
but the sum is not direct because $U\cap W=U\neq\{(0,0)\}$. Equivalently,
the vector $(1,0)$ has two different decompositions:
\[
(1,0)=(1,0)+(0,0)=(0,0)+(1,0),
\]
where the first term in each sum belongs to $U$ and the second belongs to
$W$.
\end{example_box}

\textbf{External Direct Sum} Let ${(V_i)}_{i\in I}$ be a family
of vector spaces over the same field $\mathbb{F}$. Their \emph{external direct sum} is
the vector space
\[
\bigoplus_{i\in I}V_i
=\left\{{(v_i)}_{i\in I}:v_i\in V_i
\text{ and }v_i=0\text{ for all but finitely many }i\right\},
\]
with componentwise addition and scalar multiplication. Its elements are
tuples, not sums taken inside a common ambient vector space.

\textbf{Internal versus External Direct Sum} Let
\[
U=\operatorname{span}((1,0))
\qquad\text{and}\qquad
W=\operatorname{span}((0,1))
\]
be the coordinate axes in $V=\mathbb{R}^2$.

As an internal direct sum,
\[
\mathbb{R}^2=U\oplus W.
\]
Indeed, $U+W=\mathbb{R}^2$ and $U\cap W=\{(0,0)\}$. An element of this
internal direct sum is an ordinary vector $(a,b)\in\mathbb{R}^2$, and its
unique decomposition is
\[
(a,b)=(a,0)+(0,b),
\qquad (a,0)\in U,\quad (0,b)\in W.
\]
Thus, the summands are subspaces of the same ambient vector space, and their
vectors are added inside that space.

As an external direct sum, however, we construct the new vector space
\[
U\oplus_{\mathrm{ext}}W
=\{((a,0),(0,b)):a,b\in\mathbb{R}\}.
\]
Its elements are ordered pairs of vectors. For example,
\[
((2,0),(0,3))\in U\oplus_{\mathrm{ext}}W,
\]
whereas the corresponding element of the internal direct sum is the single
vector
\[
(2,0)+(0,3)=(2,3)\in\mathbb{R}^2.
\]

\textbf{Direct Product} Let ${(V_i)}_{i\in I}$ be a family of
vector spaces over the same field $\mathbb{F}$. Their \emph{direct product} is
\[
\prod_{i\in I}V_i
=\{{(v_i)}_{i\in I}:v_i\in V_i\text{ for every }i\in I\}.
\]
Addition and scalar multiplication are defined componentwise:
\[
{(v_i)}_{i\in I}+{(w_i)}_{i\in I}={(v_i+w_i)}_{i\in I},
\qquad
a*{(v_i)}_{i\in I}={(a*v_i)}_{i\in I}.
\]
With these operations, $\prod_{i\in I}V_i$ is a vector space over $\mathbb{F}$.

If $I=\{1,\ldots,m\}$ is finite, this is written
\[
V_1\times\cdots\times V_m,
\]
and the direct product and external direct sum contain the same tuples. If
$I$ is infinite, the direct product permits infinitely many nonzero
components, whereas the direct sum
\[
\bigoplus_{i\in I}V_i
\]
contains only tuples with finitely many nonzero components. Thus,
\[
\bigoplus_{i\in I}V_i\subseteq\prod_{i\in I}V_i.
\]


\section{Bonus: Linear Algebra and Set Theory}
\textit{Read only if you have read the Set Theory portions involved.}

\textbf{Theorem (all infinite bases have the same cardinality)} Let $B$ and
$C$ be two bases of an infinite-dimensional vector space $V$. Assuming the
usual axiom of choice, we have
\[
|B|=|C|.
\]

\textit{Proof:} For each $b\in B$, write $b$ in its unique basis expansion
with respect to $C$. Because a basis expansion is a finite linear combination,
there is a finite subset $C_b\subseteq C$ such that
\[
b\in\operatorname{span}(C_b).
\]
Let
\[
C_0=\bigcup_{b\in B}C_b.
\]
Every element of $B$ lies in $\operatorname{span}(C_0)$. Since $B$ spans
$V$, it follows that $C_0$ also spans $V$.

We claim that $C_0=C$. Certainly $C_0\subseteq C$. If there were some
$c\in C\setminus C_0$, then, because $C_0$ spans $V$, the vector $c$ would
belong to $\operatorname{span}(C_0)$. Thus $c$ would be a linear combination
of other elements of the basis $C$, contradicting the linear independence of
$C$. Hence $C_0=C$.

Therefore $C$ is a union of $|B|$ finite sets. Since $B$ is infinite,
cardinal arithmetic gives
\[
|C|\leq |B|\cdot\aleph_0=|B|.
\]
Interchanging the roles of $B$ and $C$ gives
\[
|B|\leq |C|.
\]
By the Cantor-Schr\"oder-Bernstein theorem, these two inequalities imply
\[
|B|=|C|.
\]
Thus every basis of $V$ has the same cardinality. \hfill$\square$

\begin{example_box}
\textbf{Example (a nonconstructively obtained basis).} Regard $\mathbb{R}$ as
a vector space over $\mathbb{Q}$. A basis of this vector space is called a
\emph{Hamel basis of $\mathbb{R}$ over $\mathbb{Q}$}. Thus a Hamel basis is a
subset $B\subseteq\mathbb{R}$, a set of real numbers $\{u_\alpha\}$ such that every real number $\beta$ has a unique representation of the form
\[
\beta = r_1 u_{\alpha_1} + r_2 u_{\alpha_2} + \cdots + r_n u_{\alpha_n},
\]
where all $u_{\alpha_i}$ are unique and $r_i\in\mathbb{Q}\setminus\{0\}$ and $n$ depends on $\beta$.
The existence of such a basis follows from Zorn's lemma, and hence from the
axiom of choice, but the proof does not produce a concrete list or formula
for all its elements. Moreover, every Hamel basis of $\mathbb{R}$ over
$\mathbb{Q}$ is uncountable: if $B$ were countable, then the set of all finite
rational linear combinations of elements of $B$ would be countable, whereas
$\mathbb{R}$ is uncountable. 
\end{example_box}

\textbf{Theorem (every vector space has a basis)} Assuming the axiom of
choice, every vector space $U$ has a basis.

\textit{Proof:} Let
\[
\mathcal{I}=\{L\subseteq U:L\text{ is linearly independent}\},
\]
ordered by inclusion. The collection $\mathcal{I}$ is nonempty because the
empty set is linearly independent.

Let $\mathcal{C}$ be a chain in $\mathcal{I}$ and define
\[
L_{\mathcal{C}}=\bigcup_{L\in\mathcal{C}}L.
\]
We claim that $L_{\mathcal{C}}$ is linearly independent. Any linear relation
uses only finitely many vectors from $L_{\mathcal{C}}$. Because
$\mathcal{C}$ is totally ordered by inclusion, all these finitely many
vectors lie together in some member $L\in\mathcal{C}$. Since $L$ is linearly
independent, all coefficients in the relation must be zero. Therefore
$L_{\mathcal{C}}\in\mathcal{I}$, and it is an upper bound for the chain.

Every chain in $\mathcal{I}$ thus has an upper bound. By Zorn's lemma,
$\mathcal{I}$ has a maximal element; call it $B_U$. It remains to show that
$B_U$ spans $U$. If some $u\in U$ did not belong to
$\operatorname{span}(B_U)$, then
\[
B_U\cup\{u\}
\]
would still be linearly independent, contradicting the maximality of $B_U$.
Hence $\operatorname{span}(B_U)=U$. Thus $B_U$ is both linearly independent
and spanning, so it is a basis of $U$. \hfill$\square$

\textbf{Theorem (complements in infinite-dimensional vector spaces)} Let
$V$ be an infinite-dimensional vector space over $\mathbb{F}$, and let $U$ be any
subspace of $V$. Assuming the axiom of choice, there exists a subspace
$W\subseteq V$ such that
\[
V=U\oplus W.
\]

\textit{Proof:} By Zorn's lemma, $U$ has a basis $B_U$. Consider the
collection
\[
\mathcal{L}
=\{L\subseteq V:B_U\subseteq L\text{ and }L\text{ is linearly independent}\},
\]
ordered by inclusion. This collection is nonempty because $B_U\in\mathcal{L}$.

Let $\mathcal{C}$ be a chain in $\mathcal{L}$ and set
\[
L_{\mathcal{C}}=\bigcup_{L\in\mathcal{C}}L.
\]
Every finite subset of $L_{\mathcal{C}}$ is contained in one member of the
chain $\mathcal{C}$ and is therefore linearly independent. Hence
$L_{\mathcal{C}}$ is linearly independent, contains $B_U$, and is an upper
bound for $\mathcal{C}$ in $\mathcal{L}$. By Zorn's lemma, $\mathcal{L}$ has
a maximal element $B$.

The set $B$ spans $V$. Indeed, if some $v\in V$ did not belong to
$\operatorname{span}(B)$, then $B\cup\{v\}$ would be linearly independent,
contradicting the maximality of $B$. Thus $B$ is a basis of $V$ extending
$B_U$.

Let
\[
C=B\setminus B_U
\qquad\text{and}\qquad
W=\operatorname{span}(C).
\]
Because $B=B_U\cup C$ spans $V$, every vector in $V$ is a sum of a vector in
$U=\operatorname{span}(B_U)$ and a vector in $W$. Therefore $U+W=V$.

Now suppose $x\in U\cap W$. Then $x$ has one basis expansion using elements
of $B_U$ and another using elements of $C$. Subtracting these expansions
gives a finite linear combination of elements of the linearly independent set
$B_U\cup C=B$ equal to $0_V$. All coefficients must therefore be zero, so
$x=0_V$. Hence
\[
U\cap W=\{0_V\}.
\]
Consequently, $V=U\oplus W$, and $W$ is a complementary subspace of $U$.
\hfill$\square$

\end{document}
